\chapter*{Resumen}

\section*{\tituloPortadaVal}

El estudio de grandes corpus musicales constituye una herramienta fundamental en disciplinas como la musicología histórica, la etnomusicología y el análisis musical. Tradicionalmente, estas investigaciones han dependido de procesos manuales de lectura y catalogación de partituras, una metodología rigurosa pero con severas limitaciones en términos de escalabilidad frente al creciente volumen de material musical digitalizado. A pesar de la aparición de formatos simbólicos estandarizados como MusicXML o MEI, persisten desafíos críticos relacionados con la heterogeneidad de los corpus, la inconsistencia en los estándares simbólicos, la complejidad algorítmica para detectar fenómenos rítmicos y la división entre el análisis estructural y semántico de las letras. Asimismo, los modelos de lenguaje (LLM) de propósito general presentan serias limitaciones de contexto y razonamiento para ejecutar análisis musicales especializados en tiempo real sobre colecciones extensas.

En este contexto, este trabajo presenta el diseño e implementación de una arquitectura automatizada, robusta y escalable orientada al análisis de corpus musicales extensos y heterogéneos. El sistema propuesto supera las limitaciones actuales mediante una infraestructura intermedia que combina el procesamiento de partituras, la extracción y normalización automática de metadatos y su guardado posterior en una base de datos SQLite, y la integración del análisis musical y semántico, llegando a adelantar en diversas tareas musicológicas especializadas a soluciones basadas en modelos comerciales de última generación como Gemini 2.5 Flash cuando estos operan únicamente mediante mecanismos básicos de recuperación documental (RAG).

\section*{Palabras clave}
   
\noindent Musicología computacional, Corpus musicales, Extracción de características, Modelos de lenguaje (LLM), Music21, Música folclórica, Recuperación de información, Model Context Protocol (MCP), Text-to-SQL, Humanidades digitales